% =============================================================
%  Section 2 -- Institutional Background and Related Literature
%  Label: sec:background
% =============================================================

\section{Institutional Background and Related Literature}%
\label{sec:background}

% -------------------------------------------------------------
\subsection{Food Prices in Chinese Inflation}\label{sec:food_cpi}
% -------------------------------------------------------------

Food is the single largest category in the Chinese CPI basket.
The National Bureau of Statistics weights food and non-alcoholic
beverages at approximately 28 percent of the overall index,
roughly twice the comparable share in high-income economies.
This weight reflects both the expenditure share of food in
Chinese household budgets and the NBS's plutocratic aggregation
method.
Because food expenditure shares decline with income (Engel's
law), the effective food weight has been falling over time---from
roughly 34 percent in the early 2000s to 28 percent after the
2016 basket revision---but remains substantially higher than
in OECD economies.

The NBS constructs the CPI using a modified Laspeyres index
with a five-year base-period update cycle.
A consequential change occurred in 2016, when the NBS merged
the previously separate urban and rural CPI baskets into a
unified national basket.
Before 2016, rural and urban price indices used different
item lists and different outlet samples, creating
discontinuities in province-level CPI series at the merge
point.
I begin the analysis in 2016 to work with a consistent
basket throughout.

Within food, pork alone accounts for around 60 percent of total
meat consumption in China.
The NBS does not publish the exact weight of individual food
items, but researchers who have backed out pork's contribution
by regressing headline CPI on sub-indices consistently place the
pork weight at 2.5--4.5 percent for the 2016--2020 basket, with
higher effective weights in provinces where per-capita pork
consumption exceeds the national mean.

That official weight understates pork's effective salience.
Pork is purchased frequently: unlike durable categories whose
prices households observe intermittently, pork is a daily
staple for a large share of the population.
Purchase frequency amplifies the informational weight that
consumers assign to price changes relative to budget share
\citep{Cavallo2017}.
The same logic applies with particular force in China, where
wet markets are the primary retail channel for fresh protein
and where household members observe pork prices at every
market visit---often daily in urban areas.

Historically, pork price cycles have been the single most
important contributor to headline CPI volatility in China.
The so-called ``pig cycle'' has generated four distinct peaks
in pork retail prices since 2005: in 2007--2008, 2010--2011,
2015--2016, and the ASF-driven peak of 2019--2020.
During each episode, year-on-year pork price inflation exceeded
50 percent for at least two consecutive months, and the pork
sub-index accounted for 60--80 percent of the total variation
in food CPI.
The 2019 episode was the most extreme in the modern CPI
record: the NBS meat and poultry sub-index reached 122
percent year-on-year in February 2020.

The baseline 28 percent food weight creates a strong prior
that food-price movements are informative about aggregate
inflation.
When pork is expensive, households that weight recent
experience heavily form upward-revised expectations even if
non-food prices are stable.
This channel is not irrational in a narrow accounting sense
because food prices enter headline CPI.
The issue is magnitude.
The expectation response can exceed what food weights justify.

% -------------------------------------------------------------
\subsection{African Swine Fever as a Natural Experiment}%
\label{sec:asf}
% -------------------------------------------------------------

African Swine Fever (ASF) is a highly contagious viral disease
affecting domestic and wild pigs with near-100 percent mortality
in infected herds.
It has no licensed vaccine, no cure, and no transmission pathway
to humans.
These features make ASF exceptionally difficult to contain
once it enters a country's swine population.

China confirmed its first domestic case in Shenyang, Liaoning
province, in August 2018.
Epidemiological tracing suggested that contaminated swill feed
and long-distance live-pig transport were the initial vectors.
Within two months, ASF had been confirmed in Heilongjiang,
Henan, Jiangsu, and Zhejiang, spanning a 2,000-kilometer
north-south corridor.
By February 2019, every mainland province except Hainan had
reported at least one outbreak.

The geographic pattern of spread reflected China's live-pig
trade network.
Northeastern provinces (Liaoning, Heilongjiang, Jilin) are
net exporters of pigs to consuming regions in the southeast
(Guangdong, Fujian, Zhejiang).
Long-distance truck transport moved infected animals across
provincial borders before local outbreaks were detected.
The Ministry of Agriculture imposed interprovincial transport
bans in September 2018, but enforcement was incomplete, and
small-scale farmers---roughly 45 percent of national pork
production---continued informal cross-border sales.

Between August 2018 and mid-2019, the official culling count
exceeded 1.1 million pigs, but independent estimates from
the FAO placed the total herd reduction at 30 to 55 percent
of the pre-ASF national stock---roughly 120 to 210 million
animals.
Wholesale pork prices remained roughly flat through the first
five months after the initial outbreak as frozen stocks buffered
supply.
Starting in March 2019, prices began a sustained increase.
By October 2019, the NBS retail pork price index stood at
approximately double its August 2018 level.
The price increase was geographically uneven: provinces with
larger pre-ASF herd densities (Sichuan, Henan, Hunan, Hubei)
experienced the sharpest increases, while coastal provinces
with cold-chain import infrastructure partially offset
domestic supply losses with imports from the EU, Brazil,
and the US.

The price shock reversed between 2020 and 2022.
As surviving large-scale producers expanded capacity, the
national herd rebuilt faster than expected; by mid-2021, hog
inventories had returned to 95 percent of pre-ASF levels.
Pork prices fell below pre-ASF levels in late 2021.
This reversal is consequential for identification.
Households who expected continued high inflation after
observing the 2018--2020 meat-price surge were wrong about
the direction, not just the magnitude, of subsequent price
movements.

Two features of ASF matter for the identification strategy.
First, the disease spread through biological transmission in
pig populations, not as a consequence of local economic
conditions or household demand.
Cross-province variation in meat prices is therefore plausibly
orthogonal to demand shocks that might independently affect
inflation expectations.
Second, the province-level treatment intensity was determined
by pre-existing herd geography---centuries of agricultural
specialisation, climate suitability, and feed-cost
differentials---rather than by contemporary household
characteristics.
This pre-determination rules out the concern that provinces
with systematically different expectation-formation processes
selected into larger shocks.

% -------------------------------------------------------------
\subsection{Related Literature}\label{sec:literature}
% -------------------------------------------------------------

Related evidence falls into three strands.

\subsubsection{Overreaction in Expectations}

\citet{Coibion2015} provide a foundational test using the
information-rigidity framework.
Under noisy information or sticky updating, forecast revisions
should positively predict subsequent forecast errors.
A negative revision-error slope, conversely, indicates
overreaction.
\citet{Coibion2015} find that US professional forecasters
display positive revision-error covariance for most
macroeconomic variables, consistent with information rigidity
rather than overreaction.

\citet{Bordalo2020b} revisit this test with a longer panel
of professional forecasts across six macroeconomic variables.
They document a sign switch: at short horizons, forecast
revisions negatively predict errors (overreaction), while at
longer horizons the relationship is positive (underreaction).
The diagnostic expectations model of \citet{Bordalo2018}
generates exactly this horizon-dependent pattern through a
single parameter $\theta$ that governs the degree to which
agents overweight the likelihood ratio of the signal.

\citet{Angeletos2021} develop a model in which dispersed
information generates aggregate patterns observationally
similar to overreaction even when individual agents update
rationally.
The province-level design here partially addresses
this concern: because the treatment varies across provinces
within a given wave, the identification does not rely on
aggregate time-series patterns that could be confounded by
general equilibrium feedback.

I provide a household-level test that exploits
an external commodity shock rather than within-survey revisions.
The meat-shock design avoids the mechanical correlation that
arises in the standard revision-error framework: the treatment
is an NBS price statistic that does not contain any household
expectation data, and the dependent variable (forecast error)
is constructed from the household's expectation and subsequent
realised CPI---both measured independently of the treatment.

\subsubsection{Salience and Inflation Expectations}

\citet{Cavallo2017} show that US households weight frequently
purchased goods more heavily in their inflation perceptions
than the official CPI weights imply.
Respondents systematically overweight price changes for items
they buy weekly (groceries, gasoline) relative to items they
buy annually (appliances, insurance).
The overweighting is large: perceived inflation tracks a
frequency-weighted index roughly twice as closely as it tracks
the official CPI.

\citet{DacuntoHoangPirolliWeber2021} extend this logic to
European households using the European Commission's Consumer
Survey.
Country-level variation in the prices of frequently purchased
goods predicts cross-country variation in household inflation
expectations after controlling for official CPI.
The effect is asymmetric: price increases in salient categories
raise expectations by more than price decreases lower them.

\citet{Malmendier2016} document that personally experienced
inflation has a stronger influence on expectations than
historical data alone.
Using the Michigan Survey of Consumers, they show that
respondents who lived through periods of high inflation
report systematically higher expectations decades later.
\citet{Kuchler2019} document this pattern for housing
expectations, showing that personal experience with local
house-price appreciation predicts national housing return
forecasts.

These findings motivate the commodity-specific design used here.
If salience operates through purchase frequency and personal
exposure, meat and grain shocks should produce different
expectation responses despite similar CPI weights.
The placebo test uses this asymmetry: grain shocks of
comparable CPI weight do not generate forecast errors,
which rules out mechanical pass-through as the explanation
for the meat-shock result.

\subsubsection{China-Specific Evidence}

Household inflation expectations in China have received
limited attention relative to the US and European literatures.
The primary data source is the People's Bank of China (PBoC)
Urban Depositor Survey, a quarterly survey of roughly 20,000
depositors in 50 cities that uses an ordinal scale to elicit
expected price changes.

\citet{Chen2022} document how financial-policy regimes shape
China's macroeconomic dynamics.
Direct evidence on household inflation-expectation formation in
China remains limited and is often based on aggregate diffusion
indices from the PBoC Urban Depositor Survey.
The aggregate nature of these data limits identification of
cross-sectional expectation responses, and the survey samples
depositors rather than the general population.

The CFPS micro data used here address both limitations.
The CFPS is a nationally representative panel covering both
urban and rural households across 25 provinces.
Province-level treatment variation allows me to test whether
the expectation response differs across locations with
different shock magnitudes, and the panel dimension allows
me to control for household fixed effects.

The ASF episode provides identifying variation that has not
been exploited in prior work on Chinese inflation expectations.
Existing studies of ASF in the economics literature focus
on supply-chain disruption, international trade effects, and
food security, but none examines the household-level belief
response.
This analysis fills that gap by connecting province-level
price consequences of ASF to the inflation expectations of
households exposed to those prices.

